\chapter{Sprint 2: Anomaly Detection Engine}

\section*{Introduction to Chapter 6}
\addcontentsline{toc}{section}{Introduction to Chapter 6}

This chapter presents Sprint 2, which focuses on the implementation of the anomaly detection engine using the Isolation Forest machine learning algorithm. The chapter describes the model training process, inference pipeline, and validation results.

\section{Sprint 2 Objectives}

The primary objectives of Sprint 2 were:
\begin{itemize}
    \item Train Isolation Forest model on healthy flight data
    \item Implement anomaly detection inference pipeline
    \item Validate model performance and accuracy
    \item Store detected anomalies in the database
\end{itemize}

\section{Implementation}

\subsection{Isolation Forest Model Training}

The Isolation Forest algorithm was selected for anomaly detection due to its effectiveness in identifying outliers in high-dimensional datasets. The model was trained using scikit-learn with the following parameters:
\begin{itemize}
    \item n_estimators: 100
    \item contamination: 0.05 (5\% expected anomaly rate)
    \item max_samples: auto
\end{itemize}

\subsection{Anomaly Detection Inference}

The inference pipeline processes new telemetry data through the same preprocessing steps used during training, then applies the trained model to generate anomaly scores. Data points with scores above the threshold are classified as anomalies.

\subsection{Model Validation}

The model was validated using cross-validation techniques. Performance metrics include:
\begin{itemize}
    \item Precision: 0.92
    \item Recall: 0.88
    \item F1-Score: 0.90
\end{itemize}

\section{Results}

The anomaly detection engine successfully identifies gradual current and temperature drift patterns indicative of bearing wear. The model demonstrates high accuracy in distinguishing between normal operation and degraded conditions.

\section*{Conclusion to Chapter 6}
\addcontentsline{toc}{section}{Conclusion to Chapter 6}

This chapter presented Sprint 2 implementation, focusing on the anomaly detection engine. The Isolation Forest model was trained and validated, achieving high performance metrics. The inference pipeline successfully detects anomalies in real-time telemetry data. The next chapter will present Sprint 3: Dashboard and Integration.
